\section{Experiments}\label{sec:exp}

\paragraph{Protocol.}
The full protocol, including hypotheses, metrics, splits, grids and decision
rules, was committed to the repository before any result existed. Two dated
amendments, made before the corresponding results, record a change of backbone
and the addition of the simulation. Unlabelled numbers below are pre-registered;
the analysis in \S\ref{sec:posthoc} is post-hoc and labelled as such.
\emph{Data:} FracAtlas~\cite{fracatlas2023} (\FracatlasUnits\ fractured
radiographs, \FracatlasMulti\ with $\ge2$ fracture regions) and
Kvasir-SEG~\cite{kvasir2020} (\KvasirSegUnits\ polyp images, \KvasirSegMulti\
multi-polyp). Lesions are 8-connected ground-truth components labelled at native
resolution, with $\tau=0.5$. A lesion is \emph{small} if its area is below the
33rd percentile of all component areas in the dataset (a label-only definition).
\emph{Backbone:} a supervised random forest over a multiscale filter bank,
trained with 5-fold cross-fitting so that every image receives an out-of-fold
posterior. It is deliberately weak (out-of-fold Dice \FracatlasDice\ on
FracAtlas and \KvasirSegDice\ on Kvasir); see Limitations.
\emph{Splits:} \KvasirSegReps\ random tune/cal/test partitions (20/40/40\%)
per seed, with \KvasirSegSeeds\ training seeds. Baselines calibrate on tune$\cup$cal,
i.e.\ on 50\% more data than our method, which calibrates on cal only.
\emph{Metric:} the set-to-truth ratio $\sum|S|/\sum|Y|$ on test units, at a
threshold whose guarantee covers pixel FNR, lesion miss and small-lesion miss.

\paragraph{Validity.}
\ResultValidity\ (Table~\ref{tab:main}, all $\alpha$ in the supplement).
Learn-Then-Test is the most conservative method on Kvasir and at
$\alpha\le0.10$ on FracAtlas, as expected for a high-probability guarantee.

\begin{table}[t]\centering\scriptsize
\caption{Kvasir-SEG, $\alpha=0.10$, test split, mean over \KvasirSegSeeds\
seeds$\times$\KvasirSegReps\ splits. Red: a claimed risk exceeds $\alpha$
(none). $\Delta$ is the paired relative change in set size.}\label{tab:main}
\begin{tabular}{lccccr}
\toprule
Method & pixel FNR & lesion miss & small miss & set/truth & $\Delta$ vs joint-thr. \\
\midrule
Pixel CRC~\cite{angelopoulos2024crc} & 0.099 & 0.051 & 0.082 & 3.08 & -0.1\% \\
Joint, threshold & 0.098 & 0.050 & 0.081 & 3.09 & ref. \\
Joint max-form, threshold & 0.084 & 0.039 & 0.067 & 3.24 & +5.0\% \\
LTT~\cite{angelopoulos2021ltt}, threshold & 0.071 & 0.032 & 0.056 & 3.40 & +10.2\% \\
Joint, morphological~\cite{mossina2025morph} & 0.093 & 0.050 & 0.081 & 3.14 & +1.6\% \\
Joint, CCRA-S strata~\cite{luo2025ccra} & 0.092 & 0.037 & 0.066 & 3.06 & -0.8\% \\
\textbf{Ours}, joint & 0.096 & 0.049 & 0.080 & 3.11 & +0.7\% \\
\textbf{Ours}, max-form & 0.083 & 0.038 & 0.067 & 3.25 & +5.4\% \\
\bottomrule
\end{tabular}
\end{table}

\paragraph{H1: pixel guarantees leave multi-lesion units under-protected.}
At the pixel-CRC threshold, the lesion-miss rate averaged over all units stays
below $\alpha$ ($G<1$; Table~\ref{tab:gap}). On multi-lesion units, however, it
exceeds $\alpha$: on Kvasir-SEG, $G_{K\ge2}$ reaches 1.56 at $\alpha=0.10$
(verdict \KvasirSegHOne; FracAtlas \FracatlasHOne). The 95\% split intervals are
wide because only $\approx$29 multi-lesion units fall in each Kvasir calibration
split. This reproduces, on more data, the direction of the signal seen in a
smaller pilot.

\begin{table}[t]\centering\scriptsize
\caption{Gap $G=R^{\mathrm{les}}(\hat\lambda^{\mathrm{pix}})/\alpha$ on test units,
with 95\% split intervals.}\label{tab:gap}
\begin{tabular}{llccc}
\toprule
Dataset & $\alpha$ & $G$ (all) & $G\mid K{\ge}2$ & $G_{\mathrm{small}}$ \\
\midrule
FracAtlas & 0.05 & 0.64 & 0.83 [0.22, 1.64] & 0.72 [0.14, 1.46] \\
FracAtlas & 0.10 & 0.76 & 1.14 [0.59, 1.76] & 0.73 [0.33, 1.22] \\
FracAtlas & 0.15 & 0.77 & 1.26 [0.78, 1.77] & 0.81 [0.44, 1.23] \\
\midrule
Kvasir-SEG & 0.05 & 0.40 & 0.95 [0.00, 2.24] & 0.73 [0.26, 1.28] \\
Kvasir-SEG & 0.10 & 0.51 & 1.56 [0.66, 2.52] & 0.82 [0.42, 1.24] \\
Kvasir-SEG & 0.15 & 0.62 & 1.49 [0.72, 2.28] & 0.81 [0.49, 1.13] \\
\bottomrule
\end{tabular}
\end{table}

\paragraph{Joint lesion-level control is nearly free and beats existing multi-risk methods.}
At $\alpha=0.10$ the pixel-only threshold is only \FracatlasPixelCrcRel\ (FracAtlas) and \KvasirSegPixelCrcRel\ (Kvasir) smaller than the joint one, and
at most 0.7\% smaller at any $\alpha$, because on these data the pixel
constraint binds. Under the same
lesion-level guarantee, morphological sets need \KvasirSegUnionDilationRel\
(Kvasir) and \FracatlasUnionDilationRel\ (FracAtlas) larger sets at
$\alpha=0.10$. Learn-Then-Test needs \KvasirSegLttThresholdRel\ and
\FracatlasLttThresholdRel. CCRA-S strata are within $\pm2\%$ on Kvasir and 7--16\%
larger on FracAtlas.

\paragraph{H2: lesion-equalized sets on real data --- not supported.}
The pre-registered efficiency claim failed at every $\alpha$ on both datasets.
Lesion-equalized sets tied joint thresholding (Kvasir \KvasirSegOursUnionRel,
FracAtlas \FracatlasOursUnionRel\ at $\alpha=0.10$; neither smaller in $\ge95\%$
of splits). This is the behaviour the theory predicts when the lesion constraint
does not bind. There is then nothing for equalization to save, and the tuning
split often selects thresholding itself.

\paragraph{E7: when does equalization help? (controlled simulation.)}
We vary the coupling $\gamma$ between lesion size and confidence
($a\propto r^\gamma$) and the background level $b$ (Table~\ref{tab:sim}).
Lesion-equalized sets were smaller in every cell at $\gamma=0$ ($-1$ to $-4\%$),
and by up to \SimBest\ at $\gamma=\SimBestGamma$. The gain vanished at
$\gamma\ge1.5$, where the worst cell was \SimWorst. The pre-registered rule
(advantage at all $\gamma\ge1$ and increasing in $\gamma$) therefore
\SimVerdict. A post-hoc diagnostic locates the boundary. At $\gamma\ge1.5$,
76--92\% of small lesions have a peak posterior that does not exceed the image's
99th background percentile. No transform of $\hat p$ alone can then separate
them from background (Proposition~3's precondition fails), and the method falls
back to thresholding rather than doing worse.

\begin{table}[t]\centering\scriptsize
\caption{E7: change in set size (\%) of lesion-equalized vs joint thresholding,
500 splits per cell. Bold: smaller in $\ge95\%$ of splits.}\label{tab:sim}
\begin{tabular}{lccccc}
\toprule
background, $\alpha$ & $\gamma=0$ & $\gamma=0.5$ & $\gamma=1$ & $\gamma=1.5$ & $\gamma=2$ \\
\midrule
$b=0.05$, $\alpha=0.05$ & \textbf{-4.0} & \textbf{-25.7} & -19.8 & +0.7 & +0.8 \\
$b=0.05$, $\alpha=0.10$ & \textbf{-2.1} & -41.8 & \textbf{-30.7} & -0.5 & -0.5 \\
$b=0.05$, $\alpha=0.15$ & \textbf{-1.3} & \textbf{-40.0} & \textbf{-45.0} & -0.7 & +0.3 \\
$b=0.15$, $\alpha=0.05$ & \textbf{-3.6} & \textbf{-33.5} & -0.2 & +0.7 & +1.8 \\
$b=0.15$, $\alpha=0.10$ & \textbf{-2.6} & \textbf{-47.2} & -2.1 & +1.2 & +0.9 \\
$b=0.15$, $\alpha=0.15$ & \textbf{-2.0} & -38.1 & +4.0 & +1.3 & +0.8 \\
\bottomrule
\end{tabular}
\end{table}

\paragraph{Post-hoc: a conditional guarantee for multi-lesion units.}\label{sec:posthoc}
Because $L^{\mathrm{les}}$ restricted to units with $K\ge2$ is defined on a
label-determined subpopulation, it can be controlled exactly like
$L^{\mathrm{small}}$. Adding it halves the multi-lesion miss rate on Kvasir
(Table~\ref{tab:e8}), at a price of 16--28\% larger sets.

\begin{table}[t]\centering\scriptsize
\caption{Post-hoc (not pre-registered): adding the conditional risk
$L^{\mathrm{les}}\mid K\ge2$ to the joint threshold.}\label{tab:e8}
\begin{tabular}{llcccc}
\toprule
Dataset & $\alpha$ & miss$\mid K{\ge}2$, joint & miss$\mid K{\ge}2$, +cond. & set/truth, joint & set/truth, +cond. \\
\midrule
FracAtlas & 0.05 & 0.040 & 0.038 & 37.14 & 37.93 (+2\%) \\
FracAtlas & 0.10 & 0.113 & 0.092 & 27.85 & 29.89 (+7\%) \\
FracAtlas & 0.15 & 0.187 & 0.143 & 22.68 & 25.70 (+13\%) \\
Kvasir-SEG & 0.05 & 0.048 & 0.024 & 3.74 & 4.80 (+28\%) \\
Kvasir-SEG & 0.10 & 0.155 & 0.076 & 3.09 & 3.59 (+16\%) \\
Kvasir-SEG & 0.15 & 0.224 & 0.125 & 2.64 & 3.26 (+24\%) \\
\bottomrule
\end{tabular}
\end{table}

\paragraph{Limitations.}
(i)~The backbone is weak, and FracAtlas results describe a near-failing
segmenter (Dice \FracatlasDice). A GPU U-Net and the 3-D MSD Task08 pipeline are
implemented but were not run for this draft, so H2 has not been tested in the
regime where it is most plausible: a strong model whose remaining misses are
faint small lesions.
(ii)~Cross-fitted posteriors come from five models, so exchangeability is
approximate.
(iii)~Neither dataset has patient identifiers, so splits are image-wise.
(iv)~The multi-lesion subpopulation is small (29--95 units per calibration split).
(v)~The guarantees are marginal and conditional on lesion presence. None of this
is clinically validated.

\section{Conclusion}
Pixel-level conformal guarantees are carried by large lesions. On real data they
leave multi-lesion images missing lesions at up to 1.56$\alpha$. Lesion-level
risks can be controlled jointly with pixel FNR at almost no cost, and more
efficiently than with Learn-Then-Test or morphological sets. The proposed
lesion-equalized sets give large savings exactly when small lesions are present
but under-confident, and degrade gracefully to thresholding otherwise. Whether
that regime is common for strong segmenters is the open empirical question this
release is built to answer.
